\documentclass[twocolumn,10pt,journal]{IEEEtran}

\usepackage[utf8]{inputenc}
\usepackage[spanish]{babel}
\usepackage{amsmath}
\usepackage{amssymb}
\usepackage{graphicx}
\usepackage{booktabs}
\usepackage{url}
\usepackage{cite}
\usepackage{multirow}

\begin{document}

\title{Optimizaci\'on de Estructuras de Indexaci\'on Aprendidas mediante Radix Enrutamiento, B\'usqueda Branchless y Landmark Traversal}

\author{
    Elmer Valerio G\'omez Alcos, David Fern\'andez Chambilla,\\
    Joel Dandy Tintaya Cahuapaza, Luis Fernando Talizo Chambilla\\
    \small{Escuela de Posgrado, Doctorado en Ciencias de la Computaci\'on}\\
    \small{Facultad de Ingenier\'ia Estad\'istica e Inform\'atica}\\
    \small{Universidad Nacional del Altiplano}\\
    \small{Puno, Per\'u}\\
    \small{Email: \{valesitoo, davi97.dfnc, tinta.jhoel20, luistalizo20\}@gmail.com}
}

\maketitle

\begin{abstract}
Los \'indices aprendidos (\textit{Learned Indexes}) combinan la teor\'ia de estructuras de datos con modelos predictivos de Machine Learning para estimar la posici\'on f\'isica de los registros. En este trabajo, diseñamos e implementamos una suite de optimizaci\'on de bajo nivel y hardware-aware para dos de las familias de indexaci\'on m\'as importantes en la literatura cient\'ifica: el \'Arbol B* y la Skip List. Para el \'Arbol B* + ML, implementamos la eliminaci\'on de divisiones en caliente usando multiplicadores de punto flotante pre-calculados y b\'usqueda binaria branchless (libre de saltos condicionales). Para la Skip List + ML, introducimos el enrutamiento acelerado por Landmarks binarios que reduce las operaciones de pointer-chasing en memoria de $O(\log N)$ a $O(1)$ promedio. Evaluamos estas mejoras frente a sus baselines iniciales sobre un dataset inmutable de 13.3 millones de transacciones (1.26 GB), realizando 10,000 consultas. Los resultados demuestran que el \'Arbol B* ML optimizado reduce la latencia a \textbf{0.51 $\mu$s} y la Skip List ML optimizada alcanza \textbf{2.76 $\mu$s}, logrando una mejora neta significativa del 19.4\% y 17.7\% respectivamente frente a las versiones base. Adicionalmente, el procesamiento multihilo paralelo de 8 hilos escala el rendimiento del \'Arbol B* ML a \textbf{6.35 millones de QPS}. Todas las mejoras se validaron mediante pruebas de hip\'otesis de Friedman y Wilcoxon ($p < 0.001$).
\end{abstract}

\begin{IEEEkeywords}
Learned Indexes, \'Arbol B*, Skip List, RMI, Radix, Branchless Search, Landmarks, Paralelismo Multihilo, Friedman Test.
\end{IEEEkeywords}

\section{Introducci\'on}

Los \'indices en bases de datos act\'uan como funciones de mapeo de clave a posici\'on f\'isica. Estructuras como el \'Arbol B \cite{comer1979btree} y las Skip Lists \cite{pugh1990skip} resuelven las b\'usquedas en tiempo logar\'itmico, pero conllevan un alto costo en accesos a memoria y uso de cach\'e debido a su organizaci\'on jer\'arquica \cite{kraska2018case}.

El paradigma de los \'indices aprendidos (\textit{Learned Indexes}) sustituye esta jerarqu\'ia f\'isica por una funci\'on de distribuci\'on acumulada (CDF) modelada a trav\'es de regresiones lineales \cite{kraska2018case, marcus2020benchmarking}. El modelo m\'as utilizado es el RMI (\textit{Recursive Model Index}), que entrena una estructura de modelos en etapas. Aunque en teor\'ia reduce los accesos, la emulaci\'on en software de los modelos en CPU introduce overhead en tiempo de ejecuci\'on por causa de operaciones de divisi\'on, saltos de predicci\'on fallidos y pointer-chasing \cite{ferragina2020pgm, galakatos2019fiting}.

En este art\'iculo, presentamos y validamos experimentalmente un conjunto de optimizaciones de bajo nivel para las versiones aprendidas del \'Arbol B* y la Skip List, demostrando mejoras significativas sin alterar el conjunto de datos de entrada.

\section{Metodolog\'ia y Optimizaciones}

\subsection{Optimizaci\'on del \'Arbol B* + ML (RMI)}
La implementaci\'on base de RMI realiza una divisi\'on flotante en caliente para calcular el \'indice del submodelo en la Etapa 2: $j = \lfloor \text{pred}_1 \cdot M / N \rfloor$. Dado que las divisiones toman hasta 15 ciclos de CPU, optimizamos este proceso:
\begin{enumerate}
    \item \textbf{Divisi\'on-Free RMI:} Pre-calculamos el multiplicador constante $\text{invN\textunderscore M} = M / N$ en tiempo de construcci\'on. Durante la b\'usqueda, el c\'alculo se reduce a una multiplicaci\'on flotante instant\'anea: $j = \lfloor \text{pred}_1 \cdot \text{invN\textunderscore M} \rfloor$.
    \item \textbf{B\'usqueda Local Branchless:} La b\'usqueda binaria local acotada por $\Delta$ en el arreglo ordenado suele causar fallos de predicci\'on de saltos (\textit{branch mispredictions}) en los pipelines modernos de CPU. Diseñamos un bucle de b\'usqueda branchless que utiliza la instrucci\'on condicional \texttt{CMOV} del compilador, reduciendo las latencias de ejecuci\'on por instrucci\'on.
\end{enumerate}

\subsection{Optimizaci\'on de la Skip List + ML}
La Skip List ML original calcula una regresi\'on lineal de punto flotante en cada inserci\'on, pero hereda el problema de pointer-chasing (recorrido de punteros dispersos) durante las b\'usquedas.
* \textbf{Enrutamiento Predictivo por Landmarks:} Durante la construcci\'on, indexamos un arreglo de punteros uniformes llamados *landmarks* cada $1024$ nodos en memoria contigua. En lugar de iniciar la b\'usqueda desde la cabecera en el nivel m\'aximo (p. ej., nivel 15), realizamos una b\'usqueda binaria en cach\'e sobre los landmarks. Una vez identificado el bloque correcto, el algoritmo ingresa directamente en un nivel local bajo (nivel 2 o menor), requiriendo a lo m\'as $4$ a $8$ saltos f\'isicos para resolver la b\'usqueda en $O(1)$ promedio.

\subsection{Diseño Experimental}
Evaluamos 6 estructuras (3 variantes del \'Arbol B* y 3 de la Skip List) sobre $13{,}305{,}915$ registros de transacciones financieras. La suite corre una carga de estr\'es de 10,000 consultas aleatorias y mide de forma secuencial y paralela la velocidad de construcci\'on, latencia promedio y ganancia neta.

\section{Resultados y Discusi\'on}

\subsection{Rendimiento Secuencial y Ganancia de Latencia}

La Tabla \ref{tab:secuencial} detalla los resultados comparativos secuenciales de las estructuras evaluadas.

\begin{table}[h]
\centering
\caption{Comparativa Secuencial "Antes vs. Despu\'es" (10k Consultas)}
\label{tab:secuencial}
\begin{tabular}{lccc}
\toprule
\textbf{Estructura} & \textbf{Construc.} & \textbf{Latencia} & \textbf{Ganancia} \\
 & \textbf{(ms)} & \textbf{($\mu$s)} & \textbf{vs Baseline} \\
\midrule
\'Arbol B* Tradicional & 119.17 & 1.952 & - \\
\'Arbol B* + ML (Base) & 54.06 & 0.635 & Baseline \\
\textbf{\'Arbol B* + ML (Optim.)} & \textbf{50.66} & \textbf{0.511} & \textbf{+19.4\%} \\
\midrule
Skip List Tradicional & 422.61 & 4.620 & - \\
Skip List + ML (Base) & 407.99 & 3.356 & Baseline \\
\textbf{Skip List + ML (Optim.)} & \textbf{351.89} & \textbf{2.760} & \textbf{+17.7\%} \\
\bottomrule
\end{tabular}
\end{table}

El **\'Arbol B* + ML Optimizado** redujo su latencia promedio de $0.635\ \mu$s a **$0.511\ \mu$s**, representando una aceleraci\'on neta del **19.4\%** debido al Radix/division-free enrutamiento y la b\'usqueda branchless.

Por su parte, la **Skip List + ML Optimizada** super\'o tanto a la versi\'on tradicional como a la versi\'on ML base, logrando una latencia de **$2.760\ \mu$s** (una reducci\'on de latencia del **17.7\%** frente al ML Base y del **40.2\%** frente a la Skip List Tradicional). Esto valida experimentalmente el impacto de los landmarks predictivos que evitan el tr\'ansito innecesario de punteros en los niveles m\'as altos del \'indice.

\subsection{Escalabilidad Paralela}
La paralelizaci\'on multihilo con 8 hilos del \'Arbol B* ML Optimizado redujo el tiempo de procesamiento de 10,000 consultas de $6.42$ ms a **$1.57$ ms**, logrando un speedup de **4.09x** y alcanzando una tasa de rendimiento de **6.35 millones de QPS**.

\subsection{An\'alisis de Significancia Estad\'istica}

\begin{table}[h]
\centering
\caption{Significancia Estad\'istica no Param\'etrica ($\alpha = 0.05$)}
\label{tab:stats}
\begin{tabular}{lcc}
\toprule
\textbf{Prueba Realizada} & \textbf{Estad\'istico} & \textbf{p-value} \\
\midrule
Friedman (6 grupos) & $\chi^2 = 39{,}241.49$ & $< 0.001$ \\
Wilcoxon: B* ML Base vs Opt & $W = 1.3 \times 10^7$ & $1.3 \times 10^{-113}$ \\
Wilcoxon: Skip ML Base vs Opt & $W = 2.2 \times 10^7$ & $2.5 \times 10^{-7}$ \\
\bottomrule
\end{tabular}
\end{table}

La prueba de Friedman rechaza la hip\'otesis nula ($p < 0.001$). Las pruebas Wilcoxon post-hoc confirmaron que las diferencias neta de latencia ("Antes vs. Despu\'es") son altamente significativas para ambas familias de algoritmos.

\section{Conclusiones}

Este estudio demuestra que las optimizaciones hardware-aware aplicadas sobre los modelos matem\'aticos de indexaci\'on mejoran la eficiencia en ejecuci\'on:
\begin{enumerate}
    \item La eliminaci\'on de divisiones flotantes y la b\'usqueda branchless aceleran el RMI en un 19.4\%.
    \item El enrutamiento por landmarks en la Skip List aprendida reduce los fallos de cach\'e acelerando las b\'usquedas en un 17.7\%.
    \item Todas las mejoras fueron respaldadas con pruebas de hip\'otesis estad\'isticas ($p < 0.001$).
\end{enumerate}

\bibliographystyle{IEEEtran}
\bibliography{references}

\end{document}
